\documentclass[tikz,border=10pt]{standalone}
\input{../../../docs/tikz/dag_preamble.tex}

\begin{document}

\begin{tikzpicture}[
  x=1cm,
  y=1cm,
  every node/.append style={outer sep=4pt},
  source_premise/.style={
    dag_conditional,
    draw=blue!70!black,
    fill=blue!10
  }
]

\dagPaperMetadata
  {Quantifying Spatial Under-reporting Disparities in Resident Crowdsourcing}
  {Liu, Bhandaram, Garg}
  {Nature Computational Science}
  {2024}
  {arXiv:2204.08620 v4 / DOI 10.1038/s43588-023-00572-6}

\dagPaperLegendRightOfMetadata{
\node[dag_result, dag_template_legend] (legRes) at (0,0) {Lean-checked\\paper claim};
\node[source_premise, dag_template_legend] (legModel) at (3.35,0) {Source model,\\definition, or condition};
\node[dag_unformalized, dag_template_legend] (legScope) at (6.70,0) {Outside the\\checked scope};
}
\daglegend{(legRes)(legModel)(legScope)}{Legend}

\begin{dagPaperBody}
\begin{scope}[yshift=-1.2cm]

\node[source_premise] (IncidentModel) at (0,0) {
  \textbf{Incident and first-report model} \\
  stationary Poisson incident births; duration density; \\
  nonnegative locally integrable report intensity
};

\node[source_premise] (ReportingModel) at (7,0) {
  \textbf{Homogeneous Poisson reports} \\
  positive reporting rate and exponential \\
  interarrival times
};

\node[source_premise] (StartCondition) at (14,0) {
  \textbf{Condition 1: selected start} \\
  after the first report; conditionally rate-free \\
  and independent of the future report path
};

\node[source_premise] (EndpointCondition) at (21,0) {
  \textbf{Condition 2: causal endpoint} \\
  one rate-free policy uses only visible history; \\
  window and ordered jump timeline
};

\node[source_premise] (LikelihoodModels) at (28,0) {
  \textbf{Likelihood definitions} \\
  Eq.~(2) Poisson mass; Eqs.~(4)--(5) log link; \\
  structural-zero class with $0\leq\gamma\leq1$
};

\node[dag_result] (Lemma1) at (0,-4) {
  \textbf{Lemma 1} \\
  duration-mixture detection probability; \\
  calendar-time first reports form a Poisson process, \\
  including the zero-detection case
};

\node[dag_result] (MeanDelay) at (7,-4) {
  \textbf{Homogeneous reporting delay} \\
  expected first-report delay is $1/\lambda$
};

\node[dag_result] (ShiftedProcess) at (11.5,-8) {
  \textbf{Appendix D.8.2 shifted process} \\
  after the first jump, the remaining gaps are \\
  iid exponential at the original rate
};

\node[dag_result] (Lemma2) at (17.5,-8) {
  \textbf{Lemma 2} \\
  the wait from the selected start to the next report \\
  has the exponential tail at rate $\lambda$
};

\node[dag_result] (Theorem) at (21,-4) {
  \textbf{Theorem 1 / Appendix Theorem 2} \\
  $L_\lambda=F\,p(M\mid\lambda(e-s))$ under the full \\
  selected-start and causal-endpoint conditions
};

\node[dag_result] (LikelihoodConsequences) at (28,-4) {
  \textbf{Eqs.~(3), (6), and (7)} \\
  global rate MLE; regression likelihood; \\
  zero-inflated incident and family likelihoods
};

\node[dag_result] (Prop1) at (0,-8) {
  \textbf{Proposition 1} \\
  long-run calendar-window count rate; fixed-duration-law \\
  nonidentifiability of latent incident and reporting rates
};

\node[dag_result] (CityEndpoints) at (28,-8) {
  \textbf{Eqs.~(33) and (34)} \\
  NYC and Chicago three-way minimum \\
  observation-end formulas
};

\node[dag_unformalized, text width=12.5cm] (Boundary) at (14,-12) {
  \textbf{Scope boundary} \\
  empirical fit, dataset reconstruction, simulations, causal validity of the city policies,
  and reported numerical estimates are not mathematical theorem targets
};

\draw[dag_arrow] (IncidentModel) -- (Lemma1);
\draw[dag_arrow] (Lemma1) -- (Prop1);
\draw[dag_arrow] (ReportingModel) -- (MeanDelay);
\draw[dag_arrow] (ReportingModel.east) -- ++(1.0,0) |- ([yshift=0.7cm]ShiftedProcess.north) -- (ShiftedProcess.north);
\draw[dag_arrow] (StartCondition.south) to[out=-90,in=90] (Lemma2.north);
\draw[dag_arrow] (ShiftedProcess) -- (Lemma2);
\draw[dag_arrow] (Lemma2.north east) -- (Theorem.south west);
\draw[dag_arrow] (StartCondition.south east) to[out=-45,in=180] (Theorem.west);
\draw[dag_arrow] (EndpointCondition) -- (Theorem);
\draw[dag_arrow] (Theorem) -- (LikelihoodConsequences);
\draw[dag_arrow] (LikelihoodModels) -- (LikelihoodConsequences);

\end{scope}
\end{dagPaperBody}
\end{tikzpicture}
\end{document}
